\chapter{PENDAHULUAN}
\label{chap:pendahuluan}

\section{Latar Belakang}
\label{sec:latarbelakang}

Polusi suara mendegradasi kenyamanan akustik pada berbagai lingkungan, mulai dari kabin kendaraan hingga perangkat audio portabel. \emph{Passive Noise Cancellation} (PNC) efektif meredam kebisingan frekuensi tinggi, namun performanya menurun drastis pada frekuensi rendah \parencite{chen2025,han2023}. \emph{Active Noise Cancellation} (ANC) diterapkan untuk mengatasi keterbatasan tersebut melalui prinsip interferensi destruktif secara \emph{real-time}.

Algoritma \emph{Filtered-X Least Mean Square} (FXLMS) menjadi standar dalam ANC karena mampu mengompensasi jalur sekunder dan menjaga stabilitas konvergensi \parencite{han2023,xie2022}. Idealnya, algoritma kontrol adaptif pada sistem ANC memberikan konvergensi cepat dan stabil terhadap kebisingan \emph{broadband} yang bersifat non-stasioner tanpa membebani sumber daya komputasi perangkat secara berlebihan, terutama pada perangkat tersemat berkapasitas terbatas seperti \emph{headphone}. Pada praktiknya, FXLMS standar menggunakan \emph{step-size} ($\mu$) tetap sehingga kinerjanya bergantung pada \emph{tuning} manual dan rentan terhadap variasi daya sinyal pada kebisingan \emph{broadband} tersebut \parencite{cesur2025}. \emph{Filtered-X Normalized Least Mean Square} (FXNLMS) dikembangkan untuk mengatasi kelemahan ini dengan menormalisasi \emph{step-size} secara adaptif berdasarkan daya sinyal referensi terfilter. Pada aplikasi otomotif, pendekatan ini dilaporkan meningkatkan peredaman kebisingan hingga rata-rata 5,07 dB dengan performa terbaik di bawah 1 kHz, sekaligus mempercepat konvergensi dibandingkan FXLMS standar \parencite{chen2025,cesur2025}.

Peningkatan performa tersebut disertai konsekuensi pada sisi komputasi. Normalisasi \emph{step-size} pada FXNLMS menambah beban berupa estimasi daya sinyal dan operasi pembagian pada setiap iterasi, yang berdampak pada kompleksitas implementasi di perangkat tersemat berkapasitas terbatas \parencite{xie2022,cesur2025}. Sejauh ini, penelitian yang membandingkan FXLMS dan FXNLMS umumnya hanya menyoroti satu sisi dari persoalan tersebut: performa peredaman melalui simulasi perangkat lunak \parencite{chen2025,han2023,xie2022,cesur2025}, atau efisiensi sumber daya perangkat keras pada algoritma adaptif lain secara terpisah, seperti reduksi penggunaan sumber daya logika FPGA sebesar 12,5\% hingga 53,3\% pada implementasi \emph{Signed-Error} LMS \parencite{ma_lei}. Belum ada kajian yang menempatkan kedua sisi ini secara berdampingan untuk kebisingan \emph{broadband} pada sistem ANC \emph{headphone}, sehingga belum diketahui secara empiris apakah keunggulan konvergensi FXNLMS sepadan dengan biaya komputasi tambahan yang ditimbulkannya.

Kesenjangan inilah yang mendasari tiga pertanyaan yang diajukan dalam penelitian ini. Pertama, sejauh mana FXNLMS unggul dibandingkan FXLMS standar dalam meredam kebisingan \emph{broadband} pada sistem ANC \emph{headphone}. Kedua, apakah normalisasi \emph{step-size} benar-benar mempercepat konvergensi dan menurunkan \emph{steady-state error} saat menghadapi fluktuasi sinyal \emph{broadband}, mengingat \emph{step-size} tetap pada FXLMS menghadirkan dilema klasik antara kecepatan dan kestabilan. Ketiga, seberapa besar tambahan beban komputasi yang harus ditanggung untuk memperoleh keunggulan tersebut, dilihat dari kompleksitas operasi per-iterasi maupun konsumsi sumber daya logika pada FPGA.

\section{Rumusan Masalah}
\label{sec:rumusanmasalah}

Berdasarkan latar belakang tersebut, rumusan masalah dalam tugas akhir ini adalah sebagai berikut:

\begin{enumerate}
  \item Bagaimana perbandingan kinerja algoritma FXLMS dan FXNLMS dalam meredam sinyal kebisingan tipe \emph{broadband} pada sistem ANC \emph{headphone}?
  \item Bagaimana pengaruh mekanisme normalisasi \emph{step-size} pada algoritma FXNLMS terhadap kecepatan konvergensi dan \emph{steady-state error} saat menangani fluktuasi sinyal \emph{broadband}, jika dibandingkan dengan algoritma FXLMS standar yang menggunakan \emph{step-size} tetap?
  \item Bagaimana tingkat efisiensi komputasi algoritma FXNLMS dibandingkan FXLMS standar, ditinjau dari kompleksitas operasi per-iterasi (jumlah operasi perkalian, penjumlahan, dan pembagian) serta utilisasi sumber daya logika blok DSP, MAC, dan LUT pada arsitektur FPGA, mengingat FXNLMS memerlukan operasi normalisasi tambahan?
\end{enumerate}

\section{Batasan Masalah}
\label{sec:batasanmasalah}

Batasan masalah dalam tugas akhir ini adalah sebagai berikut:

\begin{enumerate}
  \item Topologi Sistem: Perancangan arsitektur \emph{Active Noise Cancellation} (ANC) dibatasi pada sistem \emph{Feedforward Single-Channel} (SISO - \emph{Single-Input Single-Output}), yang merepresentasikan model pemrosesan akustik untuk satu sisi \emph{earpiece} pada \emph{headphone}.
  \item Cakupan Algoritma: Algoritma filter adaptif yang diimplementasikan dan dikomparasikan secara eksklusif hanya algoritma \emph{Filtered-X Least Mean Square} (FXLMS) standar dengan \emph{step-size} tetap dan \emph{Filtered-X Normalized Least Mean Square} (FXNLMS) dengan \emph{step-size} ternormalisasi.
  \item Karakteristik Sinyal Bising: Sinyal referensi atau gangguan yang dievaluasi difokuskan pada kebisingan spektrum lebar (\emph{broadband noise}) pada pita frekuensi rendah (rentang frekuensi di bawah 1000 Hz--2000 Hz). Sinyal uji direpresentasikan menggunakan \emph{white noise}, \emph{pink noise}, atau rekaman bising lingkungan nyata (\emph{real-world ambient noise}).
  \item Metrik Evaluasi Kinerja dan Efisiensi: Evaluasi kinerja dibatasi pada metrik peredaman kebisingan (\emph{noise reduction}) yang diukur dalam SNR (\emph{Signal-to-Noise Ratio}), kecepatan konvergensi, dan \emph{steady-state error} hasil simulasi. Evaluasi efisiensi komputasi dibatasi pada kompleksitas operasi aritmetika per-iterasi serta metrik sintesis logika pada platform \emph{Field-Programmable Gate Array} (FPGA), yaitu utilisasi sumber daya (LUT, \emph{Flip-Flop}, memori, dan DSP \emph{Slices}).
\end{enumerate}

\section{Tujuan}
\label{sec:tujuan}

Berdasarkan rumusan masalah, tujuan dari penelitian ini adalah sebagai berikut:

\begin{enumerate}
  \item Menganalisis secara komparatif kemampuan algoritma FXLMS dan FXNLMS dalam mereduksi sinyal kebisingan spektrum lebar pada sistem ANC \emph{headphone}.
  \item Mengukur dan mengidentifikasi pengaruh mekanisme normalisasi \emph{step-size} pada algoritma FXNLMS terhadap performa metrik kontrol, khususnya pada kecepatan konvergensi (\emph{convergence rate}) dan nilai galat kondisi tunak (\emph{steady-state error}), jika dibandingkan dengan performa FXLMS standar.
  \item Menghitung dan membandingkan tingkat efisiensi komputasi antara algoritma FXLMS dan FXNLMS, mencakup kompleksitas operasi per-iterasi serta konsumsi sumber daya perangkat keras digital LUT, \emph{Flip-Flops}, dan DSP \emph{Slices} yang dibutuhkan untuk mengimplementasikan koefisien filter \emph{broadband} pada platform FPGA.
\end{enumerate}

\section{Manfaat}
\label{sec:manfaat}

\subsection{Bagi Peneliti}
\label{sec:manfaat_peneliti}

\begin{enumerate}
  \item Meningkatkan pemahaman dan keterampilan dalam perancangan sistem \emph{Active Noise Cancellation} (ANC), khususnya algoritma filter adaptif FXLMS dan FXNLMS, aritmetika \emph{fixed-point}, serta implementasi dan sintesis HDL pada FPGA.
  \item Menerapkan pengetahuan yang diperoleh selama pendidikan sarjana pada bidang pengolahan sinyal digital dan rekayasa elektronika.
\end{enumerate}

\subsection{Bagi Masyarakat Umum}
\label{sec:manfaat_masyarakat}

\begin{enumerate}
  \item Memberikan kontribusi pengetahuan mengenai perbandingan kinerja dan efisiensi komputasi algoritma FXLMS dan FXNLMS pada sistem ANC \emph{headphone} berbasis FPGA, sebagai dasar kuantitatif dalam memilih algoritma ANC untuk perangkat tersemat berkapasitas terbatas.
\end{enumerate}